\section{Vị trí không còn nằm trong kiến trúc}
\label{sec:ch11-vi-tri-khong-duoc-cho}

\index{Vision Transformer!token lớp và vị trí}%
Self-attention của chương~\ref{ch:transformer} không biết thứ tự. Hoán vị token
đầu vào thì nó hoán vị đầu ra theo, và phép tính vẫn đúng y nguyên. Với ảnh,
điều đó còn khó chịu hơn: đổi patch ở góc trái với patch ở giữa không chỉ đổi
vị trí trong danh sách, nó đổi nghĩa của cả bố cục.

ViT thêm hai loại tham số mà ảnh không tự mang. Đầu dãy là một \tn{token lớp}{class
token}, một vector. Sau encoder, bài dùng output ở đúng vị trí ấy làm biểu diễn của cả ảnh
và đặt classifier lên nó. Bài viết `class token`; gọi nó là `CLS` trong code
thì tiện, nhưng đừng để cái tên làm nó giống một pixel. Nó không nhìn một mảnh
nào cả. Nó là một chỗ để các patch dồn thông tin về.

Thứ hai là \tn{embedding vị trí}{positional embedding} học được, một bảng có
một hàng cho mỗi vị trí. Nếu \(\vect{x}_1, \ldots, \vect{x}_N\) là các patch
đã làm phẳng và \(\vect{c}\) là class token, đầu vào encoder là

\begin{equation}
\label{eq:ch11-dau-vao-vit}
\mat{Z}_0 = [\vect{c}; \vect{x}_1\mat{E}; \ldots; \vect{x}_N\mat{E}]
  + \mat{E}_{\mathrm{pos}}.
\end{equation}

\(\mat{E}_{\mathrm{pos}}\) có một hàng cho mỗi vị trí, kể cả class token.
Nó không là mã sin-cos của chương~\ref{ch:transformer}; nó là tham số học được
từ dữ liệu. Bài thử cả 1-D, 2-D và relative position ở phụ lục D.4. Bảng của
họ không thấy lợi ích đáng kể của loại biết lưới hai chiều hơn, trong cấu hình
ablation ấy~\autocite{dosovitskiy2021vit}.

Chỗ này đáng dừng lại. Một embedding vị trí nói patch thứ bảy ở vị trí thứ bảy.
Nó không nói patch thứ bảy chỉ nên nhìn bốn patch sát nó, cũng không nói một
cạnh thẳng ở đó giống một cạnh thẳng ở patch khác. Nó đưa tọa độ vào dãy. Còn
biến tọa độ thành một quy tắc cục bộ là việc mô hình phải học.

Đó là lý do tôi không gọi ViT là \enquote{CNN viết bằng attention}. Nó có thể
học các quan hệ mà CNN buộc sẵn, như định lý Cordonnier ở
chương~\ref{ch:bias-quy-nap} đã nói. Nhưng hai chữ \enquote{có thể} đi kèm một
đánh đổi về dữ liệu.
